% !TEX root = ../../../main.tex
% 所属：第一部分：单变量分析学 / 第1章：实数理论——从Peano公理到实数的连续性
% 节：1.2 代数结构初步
% 正文按原文件顺序保存；由 ../chapter1.tex 调用。
\section{代数结构初步}


从自然数向有理数扩张的过程，是代数结构由低级到
高级的过程，这里对代数结构进行初步的说明.
\begin{definition}
	代数结构，是一个集合，配上一个或多个运算，并且满足一定的公理 .运算的一般形式有加法，乘法，数乘，甚至交并非的集合运算.
\end{definition}


\begin{definition}简单代数结构树状图
	
	
\end{definition} 
\begin{tikzpicture}[
	node distance=1cm and 0.8cm,  % 减小节点间距
	algNode/.style={
		draw=black,
		thick,
		rounded corners=4pt,     % 减小圆角
		minimum width=2.5cm,     % 减小宽度
		minimum height=0.8cm,    % 减小高度
		text width=2.3cm,        % 减小文本宽度
		align=center,
		font=\scriptsize         % 使用更小的字体
	},
	arrow/.style={
		-Stealth,
		thick,
		line width=0.6pt        % 减细箭头
	}
	]
	
	% 定义节点
	\node[algNode] (magma) {群胚\\集合 + 运算+封闭性};
	\node[algNode, below=of magma] (semi) {半群\\+ 结合律};
	\node[algNode, below=of semi] (monoid) {幺半群\\+ 单位元};
	\node[algNode, below left=of monoid] (group) {群\\+ 逆元};
	\node[algNode, below right=of monoid] (ring) {环\\加法阿贝尔群\\乘法半群或幺半群};
	\node[algNode, below=of group] (abelian) {阿贝尔群\\+ 交换律};
	\node[algNode, below=of ring] (commRing) {交换环\\乘法交换律};
	\node[algNode, below=of commRing] (domain) {整环\\无零因子+单位元};
	\node[algNode, below=of domain] (field) {域\\乘法逆元存在\\加法，乘法，分配公理};
	
	% 绘制箭头
	\draw[arrow] (magma) -- (semi);
	\draw[arrow] (semi) -- (monoid);
	\draw[arrow] (monoid) -- (group);
	\draw[arrow] (monoid) -- (ring);
	\draw[arrow] (group) -- (abelian);
	\draw[arrow] (ring) -- (commRing);
	\draw[arrow] (commRing) -- (domain);
	\draw[arrow] (domain) -- (field);
	
\end{tikzpicture}
\newpage

